\begin{table}[!h]
\centering \footnotesize
\begin{tabular}{p{.99\columnwidth}}

\vspace{-1mm}\textit{arxiv} \\ \hline
familiar with the specific, expertise in this area, research in a specialized, suggests that a significant, likely for an academic \\ 

\vspace{-1mm}\textit{PubMed} \\ \hline
knowledge about the disease, background or some familiarity, structure of a scientific,  professionals with a strong \\ 

\vspace{-1mm}\textit{SciTLDR}\\ \hline
context for these terms, specific to these fields, audience with some technical, networks and the concept, fields such as artificial
 \\ 

\vspace{-1mm}\textit{SJK} \\ \hline
easy for most adults, straightforward and the concepts, easy for an adult, readers with a basic, concepts in a clear \\ 

\vspace{-1mm}\textit{CDSR}\\\hline
 text assumes some basic, assumes some basic knowledge, basic knowledge of medical, general adult audience, medical
\\ 

\vspace{-1mm}\textit{PLOS}\\ \hline
using technical terms, require specialized knowledge, strong background knowledge, discusses complex concepts
 \\ 

\vspace{-1mm}\textit{eLife} \\ \hline
explanation of the concepts, understand for a general, explanation of these concepts, understanding of the concepts, without such a background
\\ 

\vspace{-1mm}\textit{Eureka} \\ \hline
context for a non-expert, unfamiliar to some adult, non-expert with some basic, context of the research, terms are not overly
\\ 

\vspace{-1mm}\textit{CELLS} \\ \hline
audience with no science, topic is a specialized, foundation in these fields, readers with some scientific
\\ 

\vspace{-1mm}\textit{SciNews}\\ \hline
understandable by those without, understand all the details, includes a few specialized, make it more readable, understanding of these fields \\ 
\end{tabular}
\caption{Keywords mentioned in the reasoning of the LM evaluator for why a summary was given a certain readability score, stratified by dataset.}
\label{tab:keywords_by_dataset}
\end{table}